\begin{abstract}
Computational models of the mouse brain have become a relevant issue for
neuroscience and Artificial Intelligence due to the availability of open
resources at several biological scales. These resources include large-scale
electrophysiology, visual prediction benchmarks, mesoscale anatomical
information, and local electron-microscopy reconstructions coupled with
functional recordings. However, this situation also introduces an important
methodological problem. Predictive performance, reproducibility, anatomical
plausibility, and mechanistic identifiability are often discussed together even
when they correspond to different forms of evidence. In this paper, a
reproducible framework called MouseBrainBench is presented to address this
problem. It is designed to audit partial digital models of the mouse brain
without claiming the existence of a complete mouse-brain digital twin. The
framework combines synthetic known-truth experiments, Allen Visual Behavior
Neuropixels, Sensorium static and Dynamic Sensorium benchmarks, and MICRONS
CAVE data. Several experiments have been carried out to evaluate the viability
of the proposal. The obtained results show that Allen provides reproducible but
non-identifiable evidence, Sensorium provides predictive but non-causal
evidence, and MICRONS provides the main positive empirical case. Specifically,
synaptically connected directed pairs show closer readout-location similarity
than distance- and degree-matched non-connected controls in a discovery cohort
and two independent hold-out cohorts. These results support MouseBrainBench as a
critical and reproducible audit layer for deciding which claims can be defended
by partial digital models of the mouse brain.
\end{abstract}
